\chapter{Relative value sur obligations d'État}\label{ch:s2:government-bond-relative-value}

Ajuster une courbe lisse aux obligations d'un État~: certaines se situent au-dessus, d'autres en dessous. Un book de relative value achète les obligations bon marché et vend les plus chères, couvertes,
et attend quelques semaines que la courbe les ramène à leur niveau. La courbe que la Federal Reserve Board publie pour le marché du Trésor exclut, entre autres,
les deux émissions les plus récentes à chaque échéance, car elles se traitent souvent avec une prime liée à leur liquidité et à leur valeur en repo. Sur le marché synthétique de ce chapitre,
quarante obligations présentent des erreurs de pricing de quelques points de base. Un book cher--bon marché hebdomadaire affiche un ratio de Sharpe de 3.46 avant coûts lorsque le book
est neutre en DV01, et de 5.41 lorsqu'il est neutre sur les trois facteurs de la courbe. Après un coût d'un quart de point de base par unité de DV01 échangée, ces
valeurs tombent à 1.13 et 0.69~: la meilleure couverture devient le moins bon book. Les résidus se résorbent aussi plus lentement qu'une estimation à un jour ne le laisse penser~: une demi-vie de 17.9
jours selon cette estimation, mais 40~\% d'un résidu subsistent encore après 120 jours. La construction est \texttt{firm.bondrv}.

\section{Courbes ajustées et leurs résidus}

Gürkaynak, Sack et Wright ont publié la courbe quotidienne du Trésor de la Board à partir de 1961. Elle est ajustée avec l'extension de Svensson de la forme de Nelson--Siegel,
qui décrit l'ensemble de la courbe à l'aide de quelques paramètres. Ce qui entre dans l'ajustement compte autant que la forme. La Board écarte les obligations dotées de
clauses optionnelles, les titres à moins de trois mois de l'échéance, tous les bons du Trésor (dont le marché est segmenté), les obligations à vingt ans émises à partir de 1996 (qui paraissaient souvent bon marché
face aux notes à dix ans), ainsi que les émissions on-the-run et la première off-the-run entre deux et trente ans. La courbe ajustée est une courbe pour des
obligations ordinaires, bien établies.

\begin{definition}[Résidu de la courbe ajustée]\label{def:s2:government-bond-relative-value:residual}
Un \emph{résidu de la courbe ajustée}\index{résidu de la courbe ajustée} (fitted-curve residual) est le rendement de marché d'une obligation moins le rendement qu'une courbe lisse ajustée à ses cours de marché donne à son
échéance~: positif pour une obligation bon marché par rapport à la courbe (son rendement est élevé), négatif pour une obligation chère.
\end{definition}

Trois facteurs décrivent l'essentiel des mouvements de la courbe. Sur les variations quotidiennes des rendements du Trésor à échéance constante à 1, 2, 3, 5, 7, 10 et 20 ans, de
janvier 1994 à septembre 2026, la première composante principale (un niveau) explique 87.9~\% de la variance. La deuxième (une pente) en explique 8.2~\%, et la
troisième (une courbure) 2.0~\%. \texttt{firm.bondrv} construit un marché sur cette forme. La courbe de Nelson--Siegel comporte trois facteurs à rappel vers la moyenne. Quarante obligations
démarrent avec des échéances réparties entre un et trente ans, vieillissent d'un jour à la fois et sont remplacées par une nouvelle émission à trente ans lorsqu'elles arrivent à échéance. Le
rendement de chaque obligation porte une erreur de pricing en deux parties~: 2 points de base qui reviennent vers zéro avec une demi-vie de 40 jours, et 2 autres avec une demi-vie de 250 jours. Quant
à l'analyste, il observe les rendements à travers un bruit de cotation d'un demi-point de base, comme le font les cotations périmées, et ajuste la courbe chaque jour.

\section{Butterflies et trades de courbe}

\begin{definition}[Butterfly neutre en duration]\label{def:s2:government-bond-relative-value:fly}
Un \emph{butterfly neutre en duration}\index{butterfly neutre en duration} (duration-neutral butterfly) achète (ou vend) une obligation au milieu de la courbe et vend (ou achète) des obligations de
part et d'autre, en montants dont les DV01 (livre~2, chapitre~3) se compensent et sont répartis entre les ailes de sorte que la position soit aussi neutre à une
variation de pente~; il parie sur la courbure au milieu.
\end{definition}

La mesure la plus simple est le butterfly 2--5--10 en rendements~: deux fois le rendement à cinq ans, moins ceux à deux et à dix ans. De 1994 à 2026, il a affiché en moyenne
$-1.6$ points de base, avec un écart-type de 26.6. Ses valeurs ont varié de $-79$ à 74, et sa persistance quotidienne de 0.9934 correspond à une demi-vie de 105 jours
(\cref{fig:s2:government-bond-relative-value:fly}). Un trade de courbe repose sur la même idée à deux jambes~: un steepener ou un flattener, égalisé en DV01, qui parie sur
la pente.

\begin{omfigure}
\begin{tikzpicture}
\begin{axis}[width=11cm,height=5.2cm,xlabel={\small année},ylabel={\small $2\times5$y $-$ 2y $-$ 10y (pb)},tick label style={font=\footnotesize},
  xmin=1994,xmax=2027,ymin=-80,ymax=80,grid=major,grid style={gray!25},table/col sep=comma,/pgf/number format/1000 sep={}]
\addplot[omThm,thick] table[x=year,y=fly]{figdata/strategies-2/10-government-bond-relative-value/fly.csv};
\addplot[black,thin,domain=1994:2027] {0};
\end{axis}
\end{tikzpicture}

\omcaption{Le butterfly 2--5--10 du Trésor en rendements, moyennes mensuelles des valeurs quotidiennes, de janvier 1994 à septembre 2026, à partir des
rendements à échéance constante de la Board. Calculé à partir de la série H.15 disponible sur FRED. Données~: \texttt{s2\_fetch\_ust}.}\label{fig:s2:government-bond-relative-value:fly}
\end{omfigure}

\section{Signaux cher--bon marché et leur décroissance}

\begin{definition}[Signal cher--bon marché]\label{def:s2:government-bond-relative-value:signal}
Un \emph{signal cher--bon marché}\index{signal cher--bon marché} (rich-cheap signal) classe les obligations selon leurs \omterm{def:s2:government-bond-relative-value:residual}{résidus de la courbe ajustée}, souvent normalisés par l'historique propre de chaque obligation, afin de choisir
lesquelles acheter (les moins chères) et lesquelles vendre (les plus chères) dans un book de relative value couvert.
\end{definition}

Le book (\cref{lst:s2:government-bond-relative-value:book}) se rééquilibre tous les cinq jours. Il achète les cinq obligations les moins chères et vend les cinq plus chères,
égales en DV01, et paie un quart de point de base de rendement par unité de DV01 échangée. Les positions sont ensuite couvertes de l'une de trois manières. La version
non couverte ne conserve que la jambe longue.

\begin{center}
\small
\begin{tabular}{@{}p{4.2cm}>{\centering\arraybackslash}p{2.6cm}>{\centering\arraybackslash}p{2.6cm}>{\centering\arraybackslash}p{2.6cm}>{\centering\arraybackslash}p{2.6cm}@{}}
\toprule
dix ans, DV01 brut égal à un & ratio de Sharpe (brut) & ratio de Sharpe (net) & moyenne (pb/an) & coût (pb/an)\\
\midrule
jambe longue seule & 0.36 & 0.25 & 24.4 & 7.3\\
neutre en DV01 & 3.46 & 1.13 & 11.1 & 7.3\\
neutre sur trois facteurs & 5.41 & 0.69 & 9.8 & 8.3\\
\bottomrule
\end{tabular}
\end{center}

Le book long seul est une position en duration avec une légère inclinaison, une volatilité de 67.4 points de base par an contre 3.2 pour le book neutre en DV01 et 1.8 pour
celui neutre sur les facteurs. Retirer les trois facteurs élimine les mouvements de la courbe et ne laisse que les résidus. Mais la couverture par facteurs répartit de petites positions
sur les quarante obligations, et ces positions changent chaque semaine, si bien que le turnover passe de 29 à 33 unités de DV01 par an. Au coût de 0.25 point de base, le
coût de la couverture dépasse ce qu'elle rapporte. À 0.1 point de base, le classement s'inverse~: 2.54 pour le book neutre en DV01, 3.48 pour celui neutre sur les facteurs.

La décroissance des résidus est plus lente qu'il n'y paraît (\cref{fig:s2:government-bond-relative-value:decay}). En régressant le résidu de demain sur celui d'aujourd'hui, on obtient
0.96, soit une demi-vie de 17.9 jours, plus courte que chacune des deux parties injectées, car le bruit de cotation disparaît d'un jour à l'autre. Sur des horizons plus longs, la partie lente prend
le relais. Après 20 jours, 81~\% d'un résidu subsistent encore, après 60 jours 59~\%, après 120 jours 40~\%. La seule partie rapide laisserait 71~\%, 35~\% et 12~\%.
Un book qui dimensionne sa période de détention sur l'estimation à un jour détient trop brièvement et paie un turnover excessif.

\begin{omfigure}
\begin{tikzpicture}
\begin{axis}[width=10cm,height=5.6cm,xlabel={\small horizon (jours de bourse)},ylabel={\small part du résidu restant},tick label style={font=\footnotesize},
  xmin=0,xmax=120,ymin=0,ymax=1.05,grid=major,grid style={gray!25},table/col sep=comma,
  legend style={font=\scriptsize,at={(0.5,-0.3)},anchor=north},legend columns=3]
\addplot[omThm,thick,mark=*,mark size=1.2pt] table[x=h,y=slope]{figdata/strategies-2/10-government-bond-relative-value/decay.csv};
\addplot[black,thin,dashed] table[x=h,y=planted]{figdata/strategies-2/10-government-bond-relative-value/decay.csv};
\addplot[omDef,thick,dotted] table[x=h,y=fast]{figdata/strategies-2/10-government-bond-relative-value/decay.csv};
\legend{résidus de la courbe ajustée (mesurés),erreur injectée (les deux parties),partie rapide seule}
\end{axis}
\end{tikzpicture}

\omcaption{Vitesse de décroissance des \omterm{def:s2:government-bond-relative-value:residual}{résidus de la courbe ajustée} des obligations synthétiques~: pente de la régression du résidu $h$ jours plus tard sur celui d'aujourd'hui, regroupée sur les obligations, comparée à
l'erreur injectée en deux parties et à sa seule partie rapide (demi-vie de 40 jours). Données~: \texttt{s2\_bondrv.decay}.}\label{fig:s2:government-bond-relative-value:decay}
\end{omfigure}

\section{Effet de levier et financement}

Un book dont la volatilité se chiffre à quelques points de base par an et par unité de DV01 doit être porté par un effet de levier pour compter. Les books de relative value financent leurs obligations longues en
repo (livre~2, chapitre~5) et empruntent les obligations qu'ils vendent à découvert par reverse repo. Leur rendement dépend des taux de repo propres à chaque obligation. Une obligation
qui devient special en repo est bon marché à emprunter contre du collatéral et chère à vendre à découvert. Les émissions on-the-run que la Board exclut de sa courbe sont chères en partie
à cause de cette specialness. Les risques de la stratégie viennent du financement~: décotes qui augmentent, repo qu'il est impossible de renouveler, et de nombreux books détenant les mêmes
positions. Duarte, Longstaff et Yu ont établi que les stratégies d'arbitrage sur taux qui exigent davantage de capital intellectuel ont dégagé des alphas significatifs après
facteurs obligataires et actions, et après frais. C'est aussi pour cela que les coûts de la stratégie, et non ses signaux, déterminent qui peut la mener.

\section{Fiches de stratégie}

\begin{strategyfile}{Courbe ajustée cher--bon marché}\label{strat:s2:government-bond-relative-value:richcheap}
\sfield{Qui paie, et pourquoi} Investisseurs dont la demande pour certaines émissions (rééquilibrage d'indices, détentions réglementaires, offre) les éloigne
de la courbe pendant des semaines.
\sfield{Instruments et plateformes} Obligations d'État, repo et reverse repo.
\sfield{Signal} \omterm{def:s2:government-bond-relative-value:residual}{Résidus de la courbe ajustée}, mis à l'échelle par l'historique de chaque obligation.
\sfield{Dimensionnement et exécution} Couverture en DV01 ou en facteurs de courbe~; rééquilibrage hebdomadaire ou lors de mouvements importants.
\sfield{Coûts} Fourchette sur les émissions anciennes~; taux de repo~; coût de la rotation de la couverture.
\sfield{Comment elle échoue} Résidus qui tendent (offre, specialness)~; financement qui fait défaut.
\sfield{Horizon, capacité, infrastructure} De quelques semaines à quelques mois~; un ajustement quotidien de la courbe, un desk de repo.
\sfield{Backtest rigoureux} Ajustement sur les cotations disponibles à la date~; négocier au cours acheteur et au cours vendeur~; mesurer la décroissance à l'horizon de détention.
\sfield{Sources} Gürkaynak, Sack et Wright (2007) sur la courbe ajustée~; ce chapitre~: 1.13 net neutre en DV01, 0.69 neutre aux facteurs.
\end{strategyfile}

\begin{strategyfile}{Butterfly neutre en duration}\label{strat:s2:government-bond-relative-value:fly}
\sfield{Qui paie, et pourquoi} Demande concentrée sur un point de la courbe.
\sfield{Instruments et plateformes} Trois obligations ou contrats à terme~; swaps.
\sfield{Signal} Niveau du butterfly par rapport à son historique ou à la fair value d'un modèle.
\sfield{Dimensionnement et exécution} Ailes neutres en DV01 et en pente.
\sfield{Coûts} Fourchettes des trois jambes.
\sfield{Comment elle échoue} Courbure qui tend~: la demi-vie du butterfly 2--5--10 était de 105 jours entre 1994 et 2026.
\sfield{Horizon, capacité, infrastructure} Quelques mois.
\sfield{Backtest rigoureux} Pondérations recalculées à mesure que les durations changent.
\sfield{Sources} Statistiques FRED de ce chapitre~; aucun chiffre de performance vérifié.
\end{strategyfile}

\begin{strategyfile}{Steepener ou flattener}\label{strat:s2:government-bond-relative-value:curve}
\sfield{Qui paie, et pourquoi} Personne de façon fiable~: un pari sur la pente, souvent lié à la politique monétaire.
\sfield{Instruments et plateformes} Deux obligations ou contrats à terme, appariés en DV01.
\sfield{Signal} Pente comparée à une prévision de politique monétaire ou macroéconomique.
\sfield{Dimensionnement et exécution} Neutre en DV01~; dimensionné selon la volatilité de la pente.
\sfield{Coûts} Deux jambes~; écarts de carry et de roll-down.
\sfield{Comment elle échoue} Le niveau et la pente bougent ensemble d'une façon que le ratio de couverture manque.
\sfield{Horizon, capacité, infrastructure} Quelques mois.
\sfield{Backtest rigoureux} Carry et roll-down inclus.
\sfield{Sources} Aucun chiffre de performance vérifié.
\end{strategyfile}

\begin{strategyfile}{Off-the-run contre on-the-run}\label{strat:s2:government-bond-relative-value:otr}
\sfield{Qui paie, et pourquoi} Détenteurs qui paient pour la liquidité de l'émission on-the-run et sa valeur en repo.
\sfield{Instruments et plateformes} L'émission on-the-run vendue à découvert, une émission off-the-run proche achetée.
\sfield{Signal} Le spread de rendement entre elles, comparé à son historique sur le cycle d'adjudications.
\sfield{Dimensionnement et exécution} Appariée en DV01~; la position vendeuse financée en reverse repo.
\sfield{Coûts} La specialness en repo de l'émission on-the-run~: le coût de son emprunt.
\sfield{Comment elle échoue} Une fuite vers la liquidité élargit le spread.
\sfield{Horizon, capacité, infrastructure} Un cycle d'adjudication.
\sfield{Backtest rigoureux} Taux de repo special, et non collatéral général (GC).
\sfield{Sources} Gürkaynak, Sack et Wright (2007) sur les raisons pour lesquelles les émissions on-the-run sont exclues~; aucun chiffre de performance vérifié.
\end{strategyfile}

\begin{strategyfile}{Résidu couvert par composantes principales}\label{strat:s2:government-bond-relative-value:pca}
\sfield{Qui paie, et pourquoi} Comme pour le \omterm{def:s2:government-bond-relative-value:signal}{signal cher--bon marché}.
\sfield{Instruments et plateformes} Obligations ou contrats à terme sur toute la courbe.
\sfield{Signal} Résidus après les trois premières composantes principales des variations de rendement.
\sfield{Dimensionnement et exécution} Expositions aux composantes retirées~; la couverture répartie sur de nombreuses obligations.
\sfield{Coûts} La rotation de la couverture.
\sfield{Comment elle échoue} Coûts~: dans ce chapitre, la couverture factorielle a fait baisser le ratio de Sharpe net à un quart de point de base.
\sfield{Horizon, capacité, infrastructure} Quelques semaines.
\sfield{Backtest rigoureux} Composantes estimées sur les seules données passées.
\sfield{Sources} Ce chapitre~: 5.41 brut, 0.69 net.
\end{strategyfile}

\section{Tutoriel~: au-dessus et au-dessous de la courbe}\label{tut:s2:government-bond-relative-value:tut}

\begin{tutorial}
\textbf{Objectif.} Simuler le marché des obligations d'État, ajuster la courbe chaque jour, négocier les résidus avec trois couvertures, et mesurer leur décroissance~; lire les composantes et le butterfly de la courbe
Treasury réelle. \textbf{État final~:} le tableau et les deux figures.
\begin{tutsteps}
\item \textbf{Le book}.
\omcode{code/firm/bondrv/firm_bondrv.py}{106}{131}{Positions hebdomadaires cher--bon marché, couvertes de trois façons, coûts inclus.}\label{lst:s2:government-bond-relative-value:book}
\item \textbf{La décroissance}.
\omcode{code/strategies-2/10-government-bond-relative-value/python/s2_bondrv.py}{50}{66}{Quelle part d'un résidu subsiste après h jours, par rapport aux parties injectées.}\label{lst:s2:government-bond-relative-value:decay}
\item \textbf{Exécuter} \texttt{s2\_fetch\_ust.py} une fois, puis \texttt{table()}, \texttt{table(0.1)}, \texttt{decay()} et \texttt{fig\_bondrv.py}.
\end{tutsteps}
\textbf{Pour aller plus loin.} Rééquilibrer seulement lorsqu'un résidu a bougé de plus d'un point de base~; couvrir avec trois obligations liquides au lieu des quarante~;
ajouter des émissions on-the-run dotées d'une prime de liquidité injectée, et les exclure de l'ajustement.
\end{tutorial}

\section{Construction~: relative value obligataire}\label{bld:s2:government-bond-relative-value:bondrv}

\begin{build}
\bfield{Rôle} Un marché synthétique d'obligations d'État, des ajustements quotidiens de Nelson--Siegel, les résidus, des couvertures factorielles et un book cher--bon marché.
\bfield{Interface} \texttt{BondConfig(\dots)}, \texttt{loadings(tau, decay)}, \texttt{simulate\_market(cfg)},
\texttt{fit\_residuals(yields, tau, decay)}, \texttt{neutralise(w, X)}, \texttt{book(sim, cfg, hedge)}.
\bfield{Règles} Ajuster sur les rendements que voit l'analyste~; négocier aux rendements vrais~; une obligation arrivée à échéance est clôturée à sa dernière valorisation~; coûts par unité de DV01 négociée.
\bfield{Tests d'acceptation} \texttt{code/firm/bondrv/tests/}~: les bornes des loadings~; une courbe pure ne laisse aucun résidu~; les pondérations neutralisées n'ont aucune
exposition aux facteurs~; la couverture factorielle retire les mouvements de courbe, mais pas les coûts.
\bfield{Pour aller plus loin} Le quatrième facteur de Svensson~; la specialness en repo~; les adjudications et les réémissions.
\end{build}

\begin{omsources}
\item R.~S.~Gürkaynak, B.~Sack et J.~H.~Wright, «~The U.S. Treasury yield curve: 1961 to the present~», \emph{Journal of Monetary Economics}
54(8), 2007 (document de travail FEDS 2006-28).
\item J.~Duarte, F.~A.~Longstaff et F.~Yu, «~Risk and return in fixed-income arbitrage: nickels in front of a steamroller~?~», \emph{Review of
Financial Studies} 20(3), 2007.
\item Board of Governors of the Federal Reserve System, H.15 rendements Treasury à échéance constante, via FRED.
\end{omsources}

\section{Exercices}

\begin{exercise}[$\star$]\label{exo:s2:government-bond-relative-value:1}
Une erreur de valorisation se résorbe avec une persistance quotidienne de $0.5^{1/40}$. Calculer sa demi-vie.
\end{exercise}

\begin{exercise}[$\star$]\label{exo:s2:government-bond-relative-value:2}
Le rendement à deux ans vaut 4.85~\%, celui à cinq ans 4.99~\% et celui à dix ans 5.11~\%. Calculer le butterfly 2--5--10 en points de base.
\end{exercise}

\begin{exercise}[$\star$]\label{exo:s2:government-bond-relative-value:3}
Expliquer pourquoi les résidus d'un ajustement de courbe par moindres carrés sont déjà neutres aux facteurs ajustés.
\end{exercise}

\begin{exercise}[$\star\star$]\label{exo:s2:government-bond-relative-value:4}
Pourquoi le bruit de cotation rend-il trop courte une estimation à un jour de la demi-vie des résidus~?
\end{exercise}

\begin{exercise}[$\star\star$]\label{exo:s2:government-bond-relative-value:5}
Pourquoi la courbe du Board exclut-elle les titres on-the-run, et quel trade cette exclusion suggère-t-elle~?
\end{exercise}

\begin{exercise}[$\star\star$]\label{exo:s2:government-bond-relative-value:6}
Pourquoi le livre neutre aux facteurs a-t-il un ratio de Sharpe brut plus élevé et un ratio net plus bas~?
\end{exercise}

\begin{exercise}[$\star\star\star$]\label{exo:s2:government-bond-relative-value:7}
\emph{Programmation.} Exécuter \texttt{table(0.1)}. Comment les ratios de Sharpe nets évoluent-ils, et qu'est-ce que cela dit de ceux qui peuvent mener le trade~?
\end{exercise}

\begin{exercise}[$\star\star\star$]\label{exo:s2:government-bond-relative-value:8}
\emph{Trouver la faille.} «~Nos résidus ont une demi-vie de 18 jours, donc nous clôturons chaque position au bout d'un mois.~»
\end{exercise}

\section{Problème~: au-dessus et au-dessous de la courbe}

\begin{problem}[{Problème du week-end --- un livre cher--bon marché}]\label{pb:s2:government-bond-relative-value:1}
Le marché synthétique du chapitre, la courbe du Trésor et les données publiques.

\medskip\noindent\textbf{Partie I --- La courbe.}
\begin{enumerate}
\item Que Gürkaynak, Sack et Wright ont-ils publié, et selon quelle méthode~?
\item Quels titres leur ajustement exclut-il, et pourquoi~?
\item Donner les parts de variance des trois premières composantes des variations de rendement du Trésor.
\item Définir un \omterm{def:s2:government-bond-relative-value:residual}{résidu de la courbe ajustée}.
\end{enumerate}

\medskip\noindent\textbf{Partie II --- Les trades.}
\begin{enumerate}[resume]
\item Définir un \omterm{def:s2:government-bond-relative-value:fly}{butterfly neutre en duration} et donner les statistiques du 2--5--10.
\item Définir un \omterm{def:s2:government-bond-relative-value:signal}{signal cher--bon marché}.
\item Décrire les erreurs de prix du marché synthétique.
\item Décrire le livre et ses trois couvertures.
\end{enumerate}

\medskip\noindent\textbf{Partie III --- Les résultats.}
\begin{enumerate}[resume]
\item Donner les ratios de Sharpe bruts et nets par couverture.
\item Pourquoi le livre non couvert a-t-il un ratio de Sharpe faible~?
\item Pourquoi le classement change-t-il avec les coûts~?
\item À quelle vitesse les résidus décroissent-ils~?
\end{enumerate}

\medskip\noindent\textbf{Partie IV --- Le verdict.}
\begin{enumerate}[resume]
\item Énoncer le \emph{résultat nommé}~: le ratio de Sharpe du livre cher--bon marché par couverture et la demi-vie des résidus.
\item Pourquoi la demi-vie à un jour est-elle trompeuse~?
\item Qu'ajoute l'effet de levier aux risques du trade~?
\item D'où viennent les coûts du trade~?
\item Comment mener un backtest honnête~?
\item Quelle fiche de stratégie dépend le plus du repo~?
\item Quel lien le chapitre~11 entretient-il avec ce chapitre~?
\item En une phrase~: que vend un livre d'obligations de relative value~?
\end{enumerate}
\end{problem}

\section{Questions d'entretien}

\begin{interviewq}[$\star$ \iqroles{trader}]\label{iq:s2:government-bond-relative-value:1}
Qu'est-ce qui rend une obligation chère ou bon marché~?
\end{interviewq}

\begin{interviewq}[$\star\star$ \iqroles{chercheur}]\label{iq:s2:government-bond-relative-value:2}
Comment ajusteriez-vous une courbe des taux souverains, et quelles obligations en excluriez-vous~?
\end{interviewq}

\begin{interviewq}[$\star\star$ \iqroles{trader}]\label{iq:s2:government-bond-relative-value:3}
Construisez un butterfly neutre en DV01 et en pente. Sur quoi continue-t-il de parier~?
\end{interviewq}

\begin{interviewq}[$\star\star$ \iqroles{risque}]\label{iq:s2:government-bond-relative-value:4}
Un livre de relative value dispose d'un effet de levier de 30. Qu'est-ce qui peut mal tourner~?
\end{interviewq}

\begin{interviewq}[$\star\star$ \iqroles{développeur}]\label{iq:s2:government-bond-relative-value:5}
Concevez l'ajustement quotidien de la courbe~: entrées, exclusions, contrôles et sorties.
\end{interviewq}

\begin{interviewq}[$\star\star\star$ \iqroles{chercheur}]\label{iq:s2:government-bond-relative-value:6}
Un résidu vaut $e_t + \eta_t$, avec $e_t$ un AR(1) de persistance $\phi$ et de variance $\sigma_e^2$ et $\eta_t$ un bruit indépendant de variance
$\sigma_\eta^2$. Montrez que la pente de la régression à un jour vaut $\phi\,\sigma_e^2/(\sigma_e^2 + \sigma_\eta^2)$, puis expliquez le biais de la demi-vie.
\end{interviewq}
